\documentclass[instructions.tex]{subfiles}
\begin{document}
 
\chapter{Pour être pénitent, il faut vivre dans la mortification.}
 
Ceux qui appartiennent à Jésus-Christ, dit saint Paul, ont crucifié leur chair avec ses convoitises\footnote{Qui autem sunt Christi, carnem suam crucifixerunt cum vitiis et concupiscentiis. Gal. 5, 24.}. C’est-à-dire qu’un chrétien pénitent mortifie son corps, veille sur ses sens, et sait même se mortifier dans les choses permises.
 
\primo Le corps est un ennemi dont il faut se défier. On ne doit le ménager qu’autant qu’il est nécessaire. Il est juste qu’il souffre, puisqu’il a été l’instrument du péché. Si vous l’écoutez, vous lui accorderez trop, et il ne sera jamais content. On multiplie ses besoins, quand on l’accoutume à ne se passer de rien. En lui retranchant beaucoup, il aura toujours assez.
 
La nourriture lui est nécessaire; mais il faut que la modération et la pénitence y mettent des bornes, et ces bornes sont étroites. Quelle différence y aurait-il entre vous et un animal, si vous accordiez à votre corps tout ce qu’il demande? Il vaut mieux sentir les murmures et la faim d’une chair soumise, que d’éprouver les révoltes d’un corps trop nourri et indocile. Si vous n’avez qu’une nourriture grossière ou mal apprêtée, n’en désirez pas une meilleure. Un pécheur ne mérite pas d’être traité comme un homme juste. Mangez, dit Jésus-Christ, ce qu’on vous présente, bon ou mauvais\footnote{Manducate quæ apponuntur vobis.}.
 
Si la pauvreté ne vous permet pas d’avoir le nécessaire, adorez Dieu, qui veut purifier votre âme par les besoins de votre corps. Souffrez la faim, pour imiter le jeûne de votre Sauveur. Un disciple coupable n’est pas meilleur que son maître innocent.
 
Se plaindre des jeûnes de l’Église, c’est oublier que ce sont des jours de sainteté, consacrés à une pénitence générale, des jours de salut et de grâces, qu’on doit recevoir avec joie. Si vous avez du zèle, vous vous imposerez encore, pendant l’année, d’autres jours d’abstinence. N’est-il pas juste, après avoir si souvent passé les bornes de la tempérance avec les pécheurs, de jeûner de temps en temps avec les saints? Donnez aux pauvres ce que vous vous retranchez par le jeûne; autrement, dit saint Grégoire, vous ne jeûneriez pas pour Dieu, mais pour vous, en réservant à votre corps, pour un autre jour, ce que vous lui refusez aujourd’hui. Prenez garde que l’indiscrétion et la propre volonté ne se trouvent dans vos jeûnes et dans vos mortifications. On doit mortifier le corps, mais il ne faut pas le ruiner et l’abattre. Si l’obéissance, votre santé, un directeur prudent, ne vous permettent pas de jeûner, de porter le cilice, ne le faites point.
 
\secundo Un chrétien pénitent veille sur ses sens. 1. Sur ses regards. Il fait, comme Job, un traité avec ses yeux, pour ne pas même penser à l’autre sexe. Si un objet séduisant ou une peinture indécente se présente, il dit aussitôt avec David: Ah! Seigneur, détournez mes yeux, afin qu’ils ne voient pas la vanité.
 
Les comédies et les spectacles lui paraissent indignes de l’honnête homme et du chrétien; la pompe mondaine qui les accompagne, dissipe l’esprit, séduit imperceptiblement le cœur, et fait souvent en un quart d’heure perdre plus de grâces et de vertus, que la pénitence de plusieurs années n’en a pu faire acquérir.
 
L’image du crucifix et les images des saints sont, pour un pénitent, les objets les plus agréables. Il les conserve, non comme des ornemens de sa maison, mais comme des monumens de religion, qui le font souvenir de la pénitence des saints, qu’il considère, en voyant leurs images, comme des témoins et des juges qui lui reprochent ses infidélités et son peu de courage.
 
Tout ce qu’il voit lui inspire des sentimens de pénitence. S’il regarde le ciel, il pense aussitôt, en gémissant, qu’il l’a perdu; s’il voit les arbres chargés de fruits, il se reproche à lui-même qu’il est un arbre stérile, qui ne produit que des fruits de mort; si les animaux lui obéissent et le servent, il se reproche aussitôt sa désobéissance et son ingratitude envers Dieu. Les yeux baissés, il regarde souvent la terre: Voilà, dit-il, d’où je suis sorti, et où la mort, dans peu de jours, me fera rentrer.
 
2. Il veille sur ses paroles, et parle bien plus à Dieu qu’aux créatures. Il sait que le silence préserve de bien des fautes; que la langue est comme l’instrument de tous les vices, et qu’elle tue souvent l’âme par le poison mortel dont elle est remplie\footnote{Plena veneno mortifero. Jac. 3. 8.}. Il voudrait ne se servir de la voix que pour chanter les louanges de Dieu, pour accuser ses péchés, pour prendre la défense des intérêts de Dieu et pour édifier le prochain. Il se prive même quelquefois, comme David, de parler de bonnes choses, aimant se taire par humilité\footnote{Humiliatus sum, et silui à bonis. Ps. 38.}.
 
3. Il veille sur ses oreilles, et les environne d’épines. S’il entend des paroles déréglées, il en est affligé, parce qu’il voit son Dieu offensé; s’il entend des conversations ou des lectures agréables, il ne s’y applique qu’autant qu’elles lui sont profitables. Il a soin d’écouter Dieu au fond de son âme, bien plus que les hommes. Les mondains, dit-il avec David, m’entretiennent de bagatelles; mais, Seigneur, elles ne sont pas comparables à votre loi\footnote{Narraverunt mihi iniqui fabulationes, sed non ut lex tua.}.
 
Il mortifie sa curiosité. Il sait que le plaisir d’entendre bien des choses n’est pas toujours innocent, et que, quand il le serait, le mérite est grand de s’en priver. Il regarde les concerts agréables, comme peu convenables à une ame pécheresse, qui devrait passer sa vie dans les larmes.
 
\tertio Un chrétien qui aime la pénitence, sait la pratiquer dans les choses même les plus innocentes. S’il entend une symphonie dans le lieu saint, il se réjouit de ce que tout sert à honorer le Créateur; mais en même temps il gémit de lui avoir si souvent refusé ses hommages. Il sait quelquefois se priver des plaisirs permis, pour se punir des plaisirs criminels qu’il a pris autrefois. S’il est obligé de prendre quelques récréations, loin d’y attacher son cœur, il les prend avec une secrète confusion, persuadé qu’un pécheur comme lui ne mérite que des châtimens.
 
S’il prend du sommeil, c’est à regret, se souvenant que Jésus-Christ passait les nuits en prière, et qu’à l’exemple du Fils de Dieu, au lieu d’un lit, il ne devrait pas même avoir un chevet pour reposer sa tête. Il ne prend son repos que par nécessité et parce que Dieu le veut: l’obéissance à une heure réglée l’y conduit; la paresse ne l’y retient jamais. S’il dort, ce n’est qu’autant qu’il faut pour réparer ses forces et mieux servir Dieu. S’il ne dort pas, il ne s’en inquiète point; il profite de ce temps pour pleurer ses péchés\footnote{In cubilibus vestris compungimini. Ps. 4. 5.}. S’il s’éveille, c’est pour recourir à Dieu; s’il se lève, c’est pour obéir à Dieu, et pour travailler en esprit de pénitence.
 
Oh! que de consolations ressentiriez-vous en vivant de la sorte! vous éprouveriez qu’un jour passé dans la pénitence et dans les gémissemens, est mille fois plus agréable que toutes les délices de la terre, et qu’il est bien plus doux d’expier ses péchés que de les commettre.
 
\section{Résolutions}
 
\primo Puisque mon corps a été l’instrument d’un grand nombre de péchés, je lui ferai aussi partager ma pénitence.

\secundo En conséquence, je réglerai son repos; je modérerai sa nourriture; j’observerai les jeûnes prescrits par l’Église, et je ferai souvent quelques mortifications proportionnées à mes forces et à mon travail, mais avec l’avis de mon confesseur.

\tertio Je veillerai sur mes sens, afin que la mort n’entre plus par eux dans mon âme; je mortifierai ma langue, et j’éviterai les divertissemens dangereux auxquels j’ai renoncé dans mon baptême.
 
\prayer{O mon Dieu! faites-moi la grâce de réduire mon corps en servitude, afin qu’il ne soit plus en moi l’instrument du péché, et qu’il ressuscite un jour glorieusement.}
 
\end{document}
